% ===========================================================================
\chapter{System Model and Problem Formulation}
% ===========================================================================

\section{Substrate Network Model}

The physical infrastructure is modeled as an undirected weighted graph $G_S = (N, E)$, where $N$ is the set of substrate nodes and $E$ is the set of physical links.

\subsection{Node Model}

Each substrate node $n \in N$ is characterized by a resource vector:
\begin{equation}
\mathbf{r}_n = \left(r_n^{\text{ram}},\; r_n^{\text{cpu}},\; r_n^{\text{stor}},\; c_n,\; i_n,\; a_n,\; z_n\right),
\end{equation}
where $r_n^{\text{ram}}$, $r_n^{\text{cpu}}$, $r_n^{\text{stor}}$ denote available RAM (GB), CPU (cores), and storage (GB) respectively. The three CIA security dimensions --- $c_n$ (confidentiality), $i_n$ (integrity), and $a_n$ (availability) --- are independent scores sampled from $[1, 5]$, capturing the node's security posture along each axis of the CIA triad independently. The integer $z_n \in \{0, 1, 2\}$ denotes the node's trust zone assignment: 0 (public), 1 (private), or 2 (DMZ).

The \textit{residual} resource vector $\tilde{\mathbf{r}}_n(t)$ at time $t$ reflects capacity remaining after currently active VNF allocations. CIA scores and zone assignments are static properties of the substrate node and do not change as VNFs are allocated.

\subsection{Link Model}

Each link $(u,v) \in E$ is characterized by bandwidth $b_{uv}$ (Mbps), propagation latency $\ell_{uv}$ (ms), and link security level $\sigma_{uv} \in [0.2, 1.0]$. Routing follows shortest hop-count paths; the \textit{effective latency} between nodes $u$ and $v$ is the sum of link latencies along the shortest path $P(u,v)$:
\begin{equation}
L(u,v) = \sum_{(i,j) \in P(u,v)} \ell_{ij}.
\end{equation}
The \textit{effective bandwidth} between $u$ and $v$ is the bottleneck bandwidth along $P(u,v)$:
\begin{equation}
B(u,v) = \min_{(i,j) \in P(u,v)} \tilde{b}_{ij}(t),
\end{equation}
where $\tilde{b}_{ij}(t)$ is the residual bandwidth on link $(i,j)$ at time $t$. The \textit{minimum path link security} from $u$ to $v$ is the worst-case security level along the path:
\begin{equation}
\Sigma(u,v) = \min_{(i,j) \in P(u,v)} \sigma_{ij}.
\end{equation}
This models the fact that the security level of a traffic path is determined by its weakest link.

\subsection{Topology Generation}

In the experimental configuration, the substrate network is generated as an Erd\H{o}s-R\'enyi random graph $G(N, p)$ with $|N| = 25$ nodes and edge probability $p = 1.0$ (complete graph). Node resources are sampled uniformly: RAM $\in [20, 120]$ GB, CPU $\in [10, 60]$ cores, storage $\in [250, 1500]$ GB. CIA security scores are sampled uniformly and independently per dimension: confidentiality, integrity, availability $\in [1, 5]$. Trust zones are assigned uniformly at random from $\{0, 1, 2\}$. Link bandwidths are sampled uniformly in $[100, 500]$ Mbps, latencies in $[5, 60]$ ms, and link security levels in $[0.2, 1.0]$.

\subsection{Time-to-Live and Resource Release}

Each accepted SFC placement holds resources for $\tau_r$ time steps (the SFC's TTL). Resources are released at the end of each time step via a \texttt{tick()} operation that decrements remaining TTLs and deallocates resources from expired placements. This models the dynamic, time-varying nature of real NFV deployments, where the substrate becomes gradually fragmented as active SFCs consume resources and eventually expire.

\section{SFC and VNF Request Model}

\subsection{VNF Model}

A Virtual Network Function $v$ is characterized by its resource demands:
\begin{equation}
\mathbf{d}_v = \left(d_v^{\text{ram}},\; d_v^{\text{cpu}},\; d_v^{\text{stor}}\right).
\end{equation}
VNF demands are sampled uniformly: RAM $\in [2, 12]$ GB, CPU $\in [1, 6]$ cores, storage $\in [20, 120]$ GB.

\subsection{SFC Request Model}

An SFC request $r$ is a tuple:
\begin{equation}
r = \left\langle \mathcal{V}_r,\; c_r^{\min},\; i_r^{\min},\; a_r^{\min},\; z_r^{\text{req}},\; \sigma_r^{\min},\; \ell_r^{\max},\; b_r^{\min},\; \tau_r,\; \delta_r \right\rangle,
\end{equation}
where:
\begin{itemize}
  \item $\mathcal{V}_r = (v_1, v_2, \ldots, v_{K_r})$ is the ordered sequence of $K_r$ VNFs;
  \item $c_r^{\min}, i_r^{\min}, a_r^{\min} \in [1, 5]$ are the minimum required CIA security scores at each placement node, sampled independently;
  \item $z_r^{\text{req}} \in \{0, 1, 2, -1\}$ is the required trust zone: with probability 0.3 a specific zone is required (drawn uniformly from $\{0, 1, 2\}$); otherwise $z_r^{\text{req}} = -1$ (any zone accepted);
  \item $\sigma_r^{\min} \in [0.0, 0.6]$ is the minimum link security level required on each inter-VNF path;
  \item $\ell_r^{\max} \in [35, 120]$ ms is the maximum allowed end-to-end latency;
  \item $b_r^{\min} \in [50, 250]$ Mbps is the minimum required bandwidth on each hop;
  \item $\tau_r$ is the TTL drawn from a truncated exponential distribution on $[1, 25]$;
  \item $\delta_r \in \{0,1\}$ indicates whether hard isolation is required (probability 0.1).
\end{itemize}

In our experiments, $K_r = 3$ VNFs per chain (fixed), providing a well-defined three-step decision process per SFC.

\section{Placement Constraints}

A placement $\pi_r = (n_1, n_2, \ldots, n_{K_r})$ assigns VNF $v_k$ to substrate node $n_k \in N$. A placement is \textit{feasible} if and only if all of the following constraints are satisfied:

\textbf{Resource constraints:} For each VNF $v_k$ placed on node $n_k$:
\begin{equation}
\tilde{r}_{n_k}^{\text{ram}}(t) \geq d_{v_k}^{\text{ram}}, \quad
\tilde{r}_{n_k}^{\text{cpu}}(t) \geq d_{v_k}^{\text{cpu}}, \quad
\tilde{r}_{n_k}^{\text{stor}}(t) \geq d_{v_k}^{\text{stor}}.
\end{equation}

\textbf{CIA security constraints:} For each node $n_k$ in the placement:
\begin{equation}
c_{n_k} \geq c_r^{\min}, \quad i_{n_k} \geq i_r^{\min}, \quad a_{n_k} \geq a_r^{\min}.
\end{equation}
All three CIA dimensions must be satisfied independently and simultaneously.

\textbf{Zone constraint:} If the SFC specifies a required zone ($z_r^{\text{req}} \geq 0$), all placement nodes must belong to that zone:
\begin{equation}
\forall k:\; z_{n_k} = z_r^{\text{req}}.
\end{equation}

\textbf{Bandwidth constraints:} For each consecutive pair $(n_k, n_{k+1})$:
\begin{equation}
B(n_k, n_{k+1}) \geq b_r^{\min}.
\end{equation}

\textbf{Link security constraints:} For each consecutive pair $(n_k, n_{k+1})$:
\begin{equation}
\Sigma(n_k, n_{k+1}) \geq \sigma_r^{\min}.
\end{equation}

\textbf{End-to-end latency constraint:}
\begin{equation}
\sum_{k=1}^{K_r - 1} L(n_k, n_{k+1}) \leq \ell_r^{\max}.
\end{equation}

\textbf{Hard isolation constraints:} If $\delta_r = 1$, no node in $\pi_r$ may be shared with any active SFC from a different request:
\begin{equation}
\forall k:\; n_k \notin \bigcup_{r' \neq r, r' \text{ active}} \{n_{k'}^{r'}\}.
\end{equation}
Symmetrically, nodes reserved by hard-isolated SFCs cannot be used by any other SFC.

\section{MDP Formulation}

The SFC placement problem under online request arrivals is formalized as a finite-horizon Markov Decision Process (MDP) $\mathcal{M} = \langle \mathcal{S}, \mathcal{A}, \mathcal{P}, \mathcal{R}, \gamma \rangle$.

\subsection{State Space}

The state $\mathbf{s}_t \in \mathcal{S}$ at time step $t$ encodes all information relevant to the current placement decision. It comprises:

\begin{enumerate}
  \item \textbf{Node feature matrix} $\mathbf{F} \in \mathbb{R}^{N \times 20}$: for each substrate node, a normalized vector comprising: (i)~base resources [RAM, CPU, storage, normalized by initial capacity]; (ii)~CIA security scores [confidentiality, integrity, availability, normalized by $s_{\max} = 5$]; (iii)~connectivity [average neighbor bandwidth, distance to previous placement node]; (iv)~global resource shares [RAM, CPU, storage share, VNF share, SFC tenancy]; (v)~VNF fit ratios [fit RAM, fit CPU, fit storage]; (vi)~zone one-hot encoding [is\_public, is\_private, is\_DMZ]; (vii)~link security from previous placement node.
  \item \textbf{Global context vector} $\mathbf{g} \in \mathbb{R}^{20}$: encoding current VNF demands (3), SFC CIA minima (3), SFC required zone (1), SFC minimum link security (1), SFC constraints [max latency, min bandwidth] (2), VNF progress (1), hard isolation flag (1), latency progress (1), TTL (1), remaining VNF count and aggregate demands (4), and windowed acceptance ratio (1), plus reserved padding to reach dimension 20.
  \item \textbf{Placement mask} $\mathbf{m}^{P} \in \{0,1\}^{N}$: indicating which nodes are already used in the current partial placement.
  \item \textbf{Node mask} $\mathbf{m}^{N} \in \{0,1\}^{N}$: indicating real (non-padded) nodes.
\end{enumerate}

All features are normalized to $[0, 1]$.

\subsection{Action Space}

The action space is $\mathcal{A} = \{0, 1, \ldots, N_{\max}-1\}$, where $N_{\max}$ is the maximum number of nodes (fixed at 25). At each step, only a subset $\mathcal{A}_t^{\text{valid}} \subseteq \mathcal{A}$ of feasible nodes is available, determined by the action mask $\mathbf{m}_t^{\text{act}} \in \{0,1\}^{N_{\max}}$. The agent samples from the masked softmax distribution:
\begin{equation}
\pi(\mathbf{a} \mid \mathbf{s}_t) = \text{softmax}\!\left(\mathbf{z}_t \odot \mathbf{m}_t^{\text{act}} + (1-\mathbf{m}_t^{\text{act}}) \cdot (-\infty)\right),
\end{equation}
where $\mathbf{z}_t$ is the policy network's logit vector and $\odot$ denotes element-wise multiplication.

\subsection{Transition Dynamics}

Upon action $a_t = n_k$ (placing VNF $v_k$ on node $n_k$):
\begin{enumerate}
  \item Resources $\mathbf{d}_{v_k}$ are allocated on node $n_k$: $\tilde{\mathbf{r}}_{n_k} \leftarrow \tilde{\mathbf{r}}_{n_k} - \mathbf{d}_{v_k}$.
  \item Bandwidth $b_r^{\min}$ is reserved on the path $P(n_{k-1}, n_k)$ (if $k > 1$).
  \item Latency $L(n_{k-1}, n_k)$ is added to the accumulated latency $\hat{\ell}$.
  \item The VNF index advances: $k \leftarrow k+1$.
  \item If $k = K_r$: if $\hat{\ell} \leq \ell_r^{\max}$, acceptance is triggered; otherwise, all partial resources are released and rejection is triggered.
  \item Upon SFC completion (accepted or rejected): substrate TTL is advanced via \texttt{tick()}, and the next SFC request is generated.
\end{enumerate}

\subsection{Reward Function}

The reward signal is designed to encourage high acceptance ratios while providing explicit incentives for security-aware placements:

\textbf{Intermediate step reward} (after placing VNF $v_k$, $k < K_r$):
\begin{equation}
r_t^{\text{step}} = 0.1 - 0.3 \cdot \min\!\left(\frac{L(n_{k-1}, n_k)}{\ell_r^{\max}}, 1\right) + w_{\text{link}} \cdot \max(0,\, \Sigma(n_{k-1},n_k) - \sigma_r^{\min}) + w_{\text{CIA}} \cdot \overline{m}_{\text{CIA}},
\end{equation}
where $w_{\text{link}} = 0.05$ is the link security margin weight, $w_{\text{CIA}} = 0.10$ is the CIA margin weight, and $\overline{m}_{\text{CIA}}$ is the normalized mean CIA margin of the placed node:
\begin{equation}
\overline{m}_{\text{CIA}} = \frac{1}{3 \cdot m_{\max}}\left[\max(0,\, c_{n_k} - c_r^{\min}) + \max(0,\, i_{n_k} - i_r^{\min}) + \max(0,\, a_{n_k} - a_r^{\min})\right],
\end{equation}
with $m_{\max} = s_{\max} - s_{\min} = 4$ (maximum possible margin per dimension). The $+0.1$ progress bonus rewards each feasible VNF placement; the latency term penalizes distant placements; the link security and CIA terms incentivise selecting nodes and paths with security posture exceeding the hard minimum.

\textbf{Terminal acceptance reward:}
\begin{equation}
r_T^{\text{acc}} = \left(R_{\text{acc}} + w_{\text{CIA}}^{\text{acc}} \cdot \overline{M}_{\text{CIA}}\right) \cdot \left(1 + \omega \cdot \overline{\text{AR}}_w\right),
\end{equation}
where $R_{\text{acc}} = 1$ is the base acceptance reward, $w_{\text{CIA}}^{\text{acc}} = 0.30$ is the SFC-level CIA acceptance bonus weight, $\overline{M}_{\text{CIA}}$ is the mean normalized CIA margin averaged over all $K_r$ placement nodes, $\omega = 0.5$ is the windowed AR weight, and $\overline{\text{AR}}_w$ is the acceptance ratio over the last $w = 10$ requests. The AR modifier amplifies the reward when the agent has been accepting frequently, providing curriculum-style positive reinforcement.

\textbf{Terminal rejection reward:}
\begin{equation}
r_T^{\text{rej}} = R_{\text{rej}} \cdot \left(1 + \omega \cdot (1 - \overline{\text{AR}}_w)\right),
\end{equation}
where $R_{\text{rej}} = -1$. The flipped AR modifier $(1 - \overline{\text{AR}}_w)$ amplifies the penalty when the agent has been rejecting frequently, providing stronger corrective signal during poor-performance phases.

\section{CIA Triad Security Model}

\subsection{The CIA Triad as Placement Constraints}

The CIA (Confidentiality, Integrity, Availability) triad is the foundational framework for information security assessment~\cite{editor_security_nodate}. In the context of SFC placement, each dimension captures a distinct and independently important aspect of the node's security posture:

\begin{itemize}
  \item \textbf{Confidentiality} ($c_n$): the node's ability to prevent unauthorized disclosure of data processed by hosted VNFs. A low confidentiality score indicates elevated risk of data leakage, making the node unsuitable for privacy-sensitive service chains.
  \item \textbf{Integrity} ($i_n$): the node's ability to prevent unauthorized modification of data or VNF state. A low integrity score signals higher risk of tampering, which is critical for service chains involving authentication or policy enforcement functions.
  \item \textbf{Availability} ($a_n$): the node's resilience to disruption and its capacity to remain operational under load or attack. A low availability score indicates higher risk of downtime, making the node a poor choice for latency-sensitive or mission-critical VNFs.
\end{itemize}

By treating these as three independent constraints rather than collapsing them into a single scalar, the framework enables operators to express nuanced security requirements. An SFC implementing a firewall may specify a high integrity minimum but a moderate confidentiality minimum; an SFC carrying sensitive medical traffic may specify a high confidentiality minimum but a lower availability minimum. The independent CIA constraints capture this specificity.

\subsection{Trust Zones}

Trust zones partition the substrate into security domains with different access control policies. The three zones model a standard network segmentation architecture:

\begin{itemize}
  \item \textbf{Public zone} ($z = 0$): nodes accessible from the public network; lowest trust, highest exposure.
  \item \textbf{Private zone} ($z = 1$): nodes in an internal, access-controlled segment; intermediate trust.
  \item \textbf{DMZ} ($z = 2$): nodes in a demilitarized zone between the public and private networks; used for services that must be accessible from both segments.
\end{itemize}

An SFC that specifies a required zone must have all its VNFs placed within that zone, ensuring that sensitive processing is confined to the appropriate security domain. Zone compliance is enforced as a hard constraint via action masking: candidate nodes outside the required zone receive zero weight in the action distribution.

\subsection{Link Security}

The link security attribute $\sigma_{uv} \in [0.2, 1.0]$ captures the security level of the physical or virtual link connecting nodes $u$ and $v$. This may reflect encryption strength, physical access control, monitoring intensity, or other link-level security properties. The minimum path link security $\Sigma(u, v)$ is the worst-case security level along the shortest path from $u$ to $v$, capturing the fact that the security of a traffic path is limited by its weakest link. The SFC's minimum link security requirement $\sigma_r^{\min}$ ensures that all inter-VNF paths meet a specified security floor, complementing the node-level CIA constraints with path-level security assurance.

\section{Optimization Objective}

The DRL agent seeks to maximize the expected cumulative discounted reward over an episode of $T = 300$ requests:
\begin{equation}
J(\pi) = \mathbb{E}_{\pi}\!\left[\sum_{t=0}^{T} \gamma^t r_t\right],
\end{equation}
with discount factor $\gamma = 0.99$. This objective simultaneously captures acceptance ratio (via the $R_{\text{acc}} / R_{\text{rej}}$ terms), latency efficiency (via intermediate latency penalties), CIA security posture (via the CIA margin reward terms), and link-level security (via the link security margin reward), providing a holistic measure of intelligent, security-posture-aware placement quality.
